% Propietario conservado: /Users/ruben/Documents/ChatGPT/jueces y controles/REVISION_ARGUMENTAL_APERTURAS_CIERRES_20260915/delivery/ARTICLE_V_ES_ARGUMENTAL_20260915/manuscrito/sections/40_exclusion_pauli.tex
% Recuperación documental para el artículo V; RESULTADO_RECUPERADO.
% Propietario: /Users/ruben/Documents/New project/output/REVISION_EDITORIAL_HMT_MD_20260905/01_FUENTES/INTEGRAL/tree/New project/output/ENSAMBLADO_CAUSAL_RESTAURADO_HMT_MD_20260905/fuente/manuscrito/integracion_83/parte_iv/body/B015_ch56_pauli_completacion_fermionica_7549b83967f5_04d_ocupacion_estadistica.tex
% Líneas de origen: 1--48.
% REV02: aplicación estadística; prueba común de idempotencia reunida en 20.
% Recupera además los recuentos y ocupaciones antes anticipados en 20.
% Correspondencia editorial: gestion/CONSOLIDACION_PAULI_FOCK_REV02.json.
\section{Statistical application of occupation spectra}
\label{v:v:rec:chap:ocupacion-estadistica}
\index{Pauli exclusion principle}\index{CAR}
\index{Fermi--Dirac}\index{Bose--Einstein}

The null and material realizations provide the geometric support for
occupation. The graded completion of that support determines its
multiparticle spaces and operators;
Theorem~\ref{v:v:rec:fock_gibbs:statement:01} establishes their transport under
refinement. This stage uses those results to characterize
admissible configurations, count states, and prepare the thermal evaluation.
The algebraic construction, finite counting, and equilibrium state
thus retain their respective functions and premises.

\subsection{Spectral restriction and admissible configurations}
\label{v:v:rec:pauli:section:02}

For a mode \(i\), write \(n_i=N_i=a_i^\dagger a_i\).
Theorem~\ref{v:v:rec:thm:principio-exclusion-pauli}, with its complete proof in
\eqref{v:v:pauli:eq:prueba-idempotencia}, establishes
\(n_i^2=n_i\) and \(\Spec(n_i)\subseteq\{0,1\}\). The calculation uses
only CAR and the adjoint relation, so it also applies
to any representation of those relations. Nilpotency
\((a_i^\dagger)^2=0\) prevents double fermionic occupation of the same mode.
Consequently, a fermionic configuration of \(g\) modes is a
family \((n_1,\ldots,n_g)\in\{0,1\}^g\); in the bosonic sector, the
occupations belong to \(\mathbb N_0^g\).

The fermionic representation uses the CAR algebra constructed from the
exterior completion. The relativistic spin--statistics
identification additionally retains compatibility of the rotation
representation with parity and the locality conditions stated in
Sections~\ref{v:v:pauli:car-construccion} and
\ref{v:v:rec:completaciones:section:05}. The \(4\pi\) periodicity and
exclusion are coordinated by the same sheet character, with these
representation conditions explicit.

\subsection{Dimensions of fixed-occupation sectors}
\label{v:v:rec:completaciones:section:07}

If there are \(g\) modes and \(N\) particles, the sector dimensions are
\begin{equation}
 \Omega_F(g,N)=\binom gN,
 \qquad
 \Omega_B(g,N)=\binom{g+N-1}{N}.
 \label{v:v:ocupacion:eq:dimensiones-sectores}
\end{equation}
The first formula counts subsets of modes, with zero dimension if
\(N>g\); the second counts weak compositions of \(N\) into \(g\) parts.
These counts are exact consequences of the two completions and are
obtained before introducing the thermal state.
Proposition~\ref{v:v:rec:prop:cardinalidades-microcanonicas} applies both
counts to the finite spectral screen, retaining its levels,
multiplicities, and energy and population constraints.

\subsection{Thermal readout of the graded completions}
\label{v:v:ocupacion:lectura-termica}

With a diagonal Hamiltonian \(H\), modal energy \(\epsilon\), chemical
potential \(\mu\), and thermal parameter \(\beta>0\), the mean occupations
of the two sectors are
\begin{equation}
 \bar n_{\rm par}(\epsilon)
 =\frac{1}{e^{\beta(\epsilon-\mu)}-1},
 \qquad
 \bar n_{\rm impar}(\epsilon)
 =\frac{1}{e^{\beta(\epsilon-\mu)}+1}.
 \label{v:v:ocupacion:eq:leyes-por-modo}
\end{equation}
The bosonic sector requires \(\mu<\inf\operatorname{Spec}H\).
These expressions are the per-mode specialization of
Theorem~\ref{v:v:rec:thm:factorizacion-gran-canonica-rev3}, whose proof
constructs the grand canonical trace from the admissible occupations.
The Bose--Einstein and Fermi--Dirac formulas are therefore the
thermodynamic evaluations of the symmetric and exterior completions
selected by transport. The same section derives their expressions
by varying the entropy under the energy and number constraints;
Theorem~\ref{v:v:rec:thm:limite-clasico-fd-be-rev2} proves their
common classical limit with an explicit error bound.
